\scanpage{28}

$\underline{\overline{\mathrm{XVIII}}}$ $1.2$ \mspage{12}

% Two diagonals cancel the body of the sheet.
\deleted{Pour $S$ artinien, on a $(\mathrm i)\Rightarrow(\mathrm i')$ : si $G$ vérifie $(\mathrm i)$, donc
est fini localement libre, soient $G'$ son dual de Cartier, et
$A(G)$ le schéma en groupe \inserted{(lisse)} des éléments inversibles de l'algèbre
affine de $G'$ ; $G\simeq G''$ s'identifie au sous-groupe de $A(G)$ formé
des caractères de $G'$. Soit donc une suite exacte}
\[
\deleted{0\longrightarrow G\longrightarrow G_1\longrightarrow G_2\longrightarrow0}
\]
\deleted{avec $G_1$ et $G_2$ lisses. La suite exacte longue de cohomologie
fournit, par $1.2.3$. et $1.2.2$ appliqué à $G_1$ et $G_2$}
\[
\deleted{0\longrightarrow R^1p_*G\longrightarrow\mathrm{Hom}(\mathbb G_m,G_1)\longrightarrow\mathrm{Hom}(\mathbb G_m,G_2)}
\]
\deleted{de sorte que $\mathrm{Hom}(\mathbb G_m,G)\xrightarrow{\sim}\deleted{\mathrm{Hom}} R^1p_*G$.}

\deleted{Dans le cas $(\mathrm{iii})$, le problème est local sur $S$ $(1.2.4)$.
Quitte à étendre $k$, on peut supposer l'existence d'une suite
exacte}
\[
\deleted{0\longrightarrow\mathbb G_m^n\longrightarrow G_0\longrightarrow H\longrightarrow0}
\]
\deleted{telle que $H$ ne contienne pas de tore. On applique alors
$1.2.3$ et $1.2.8$ à la suite exacte longue de cohomologie.}
% archive sheet number: 27

